% ============================================================
%  USPSC-Cap3-Metodologia_3.3-Rotulagem_Manual.tex
%  Seção 3.3 — Rotulagem Manual do Corpus
% ============================================================

\section{Rotulagem Manual do Corpus}
\label{sec:rotulagem}

A rotulagem manual constitui etapa fundamental para a construção do conjunto de referência (\textit{gold standard}) utilizado na avaliação das abordagens de classificação de sentimento descritas na \autoref{sec:abordagens}. O processo foi conduzido pelo próprio pesquisador por meio de uma ferramenta de anotação desenvolvida especificamente para este trabalho, implementada em Python com a biblioteca \texttt{streamlit} e integrada ao banco de dados PostgreSQL descrito na \autoref{subsubsec:pipeline_coleta}.

\subsection{Protocolo de Anotação em Duas Etapas}
\label{subsec:protocolo_anotacao}

O protocolo de anotação foi estruturado em duas etapas sequenciais, aplicadas a cada publicação individualmente:

\textbf{Etapa 1 — Relevância financeira.} O anotador classifica a publicação em uma de duas categorias mutuamente exclusivas: \textit{tweet financeiro} (valor~1), quando o conteúdo se relaciona diretamente ao mercado financeiro, à economia, a investimentos ou a eventos com impacto sobre ativos financeiros; ou \textit{não financeiro} (valor~0), quando o conteúdo é de outra natureza — editorial, esportivo, de entretenimento ou sem relação com o contexto econômico-financeiro. Apenas as publicações classificadas como financeiras prosseguem para a segunda etapa.

\textbf{Etapa 2 — Polaridade de sentimento.} Para as publicações identificadas como financeiras, o anotador atribui uma das três classes de polaridade:

\begin{itemize}
	\item \textbf{Positivo}: o conteúdo veicula uma perspectiva favorável ao mercado ou a agentes econômicos — crescimento, resultados acima do esperado, alívio de incertezas, aprovação de
	medidas de estímulo, entre outros;
	
	\item \textbf{Neutro}: o conteúdo tem caráter predominantemente informativo e factual, sem atribuir valoração explícita ao evento relatado;

	\item \textbf{Negativo}: o conteúdo veicula uma perspectiva desfavorável — queda de indicadores, resultados abaixo do esperado, aumento de incerteza, crises ou eventos de risco sistêmico.
\end{itemize}

A separação entre relevância financeira e polaridade de sentimento em etapas distintas é uma decisão metodológica deliberada: ela evita que publicações de conteúdo não financeiro — mesmo que textualmente negativas ou positivas — contaminem o índice de sentimento do mercado construído na \autoref{sec:indice_sentimento}.

\subsection{Ferramenta de Anotação}
\label{subsec:ferramenta_anotacao}

A ferramenta de anotação constitui o segundo estágio do \textit{pipeline} (\texttt{pipeline/02\_annotation/}) e é composta por dois arquivos:

\begin{itemize}
	\item \texttt{pipeline/02\_annotation/app.py}: \textit{launcher} multi-página Streamlit, invocado via \texttt{streamlit run pipeline/02\_annotation/app.py} (acessível em \texttt{http://localhost:8501}). Agrega em uma única interface três páginas: \textit{Classification} (\texttt{classifier.py}), \textit{Analytics} (\texttt{pipeline/03\_eda/dashboard.py}) e \textit{Exploração dos Dados} (\texttt{pipeline/03\_eda/eda.py}), permitindo ao anotador alternar entre a tarefa de rotulagem e a inspeção dos dados sem trocar de aplicação;
	
	\item \texttt{pipeline/02\_annotation/classifier.py}: página principal de anotação, que implementa a interface de classificação tweet a tweet e obtém os dados diretamente do banco de dados via
	\texttt{shared/database.py}.
\end{itemize}

A interface de classificação é composta por três componentes principais:

\begin{enumerate}
	\item \textbf{Exibição do conteúdo}: para cada publicação, a ferramenta exibe o texto integral do campo \texttt{note\_tweet} e a data de publicação, permitindo ao anotador contextualizar
	temporalmente o conteúdo antes de classificá-lo;
	
	\item \textbf{Botões de classificação}: dois botões para a decisão de relevância financeira (Etapa~1) e, condicionalmente, três botões para a atribuição de polaridade (Etapa~2). O estado da classificação corrente é destacado visualmente, permitindo revisão e correção a qualquer momento;
	
	\item \textbf{Registro de justificativa}: um formulário opcional (\textit{dialog} Streamlit) permite que o anotador registre, em texto livre, a justificativa para cada decisão — tanto para a relevância financeira quanto para a polaridade atribuída. Esse registro é armazenado nos campos \texttt{why\_is\_finance\_news} e \texttt{why\_sentiment} da tabela \texttt{tweets\_classification}, enriquecendo o \textit{gold standard} com raciocínio explícito passível de auditoria futura.
\end{enumerate}

A navegação pelo corpus é sequencial, com barra de progresso indicando o percentual de publicações já classificadas. A ferramenta prioriza automaticamente as publicações ainda sem classificação
humana, identificadas pelo campo \texttt{has\_human\_classification}. Ao final de cada publicação, um histórico tabular exibe todas as classificações anteriores registradas em \texttt{tweets\_classification}, com os campos \textit{classificador}, \textit{sentimento}, \textit{relevância financeira} e as respectivas justificativas.

\subsection{Armazenamento das Anotações}
\label{subsec:armazenamento_anotacoes}

Cada decisão de classificação é persistida em duas estruturas complementares do banco de dados:

\begin{itemize}
	\item \textbf{Tabela \texttt{tweets}}: os campos \texttt{is\_finance\_news} e \texttt{sentiment} são atualizados diretamente, tornando a classificação imediatamente disponível para as etapas de análise subsequentes;
	
	\item \textbf{Tabela \texttt{tweets\_classification}}: registra o histórico completo de cada decisão com o identificador do classificador (\texttt{classificator = ``Humano''}), as justificativas textuais e os valores atribuídos, garantindo rastreabilidade e possibilitando análises de concordância em estudos futuros.
\end{itemize}


\subsection{Particionamento do Corpus Rotulado para Treinamento e Avaliação}
\label{subsec:particionamento}

O corpus rotulado manualmente cumpre, neste trabalho, dois propósitos metodológicos distintos: (i) constitui o conjunto de referência (\textit{gold standard}) utilizado na avaliação comparativa das abordagens de classificação descritas na \autoref{sec:abordagens}; e (ii) é a fonte de exemplos supervisionados para o \textit{fine-tuning} de modelos pré-treinados ajustados localmente, como o BERTimbau (\autoref{subsec:bertimbau}). Esse duplo uso impõe uma decisão de protocolo: para que a comparação entre abordagens que viram os rótulos durante o treinamento (e.g., BERTimbau ajustado) e abordagens que não os viram (e.g., FinBERT-PT-BR e métodos léxicos) seja metodologicamente honesta, todas as abordagens devem ser avaliadas sobre o \textbf{mesmo subconjunto de teste}, definido a priori, antes de qualquer treinamento \cite{jurafsky2026,caseli2024cap1}.

Adota-se, para tanto, um particionamento estratificado em três conjuntos disjuntos --- \textbf{treinamento} (70\%), \textbf{validação} (15\%) e \textbf{teste} (15\%) ---, gerado por meio da função \texttt{train\_test\_split} da biblioteca \texttt{scikit-learn}, com semente aleatória fixa (\texttt{random\_state=42}) e estratificação pela classe de polaridade. A estratificação preserva a distribuição original das três classes (positivo, neutro, negativo) em cada conjunto, o que é particularmente relevante dado o desbalanceamento esperado no corpus em favor da classe neutra, característico de publicações jornalísticas de natureza factual \cite{loughran2011}. O particionamento é realizado uma única vez e os identificadores das publicações em cada partição são persistidos em arquivos auxiliares, garantindo que todas as abordagens avaliadas operem sobre exatamente os mesmos exemplos de teste.

A função analítica de cada partição é a seguinte:

\begin{itemize}
	\item \textbf{Treinamento (70\%)}: utilizado exclusivamente pelas abordagens que requerem aprendizado supervisionado a partir do corpus local (e.g., \textit{fine-tuning} do BERTimbau). Abordagens léxicas e o FinBERT-PT-BR não consomem esta partição;
	
	\item \textbf{Validação (15\%)}: utilizada durante o \textit{fine-tuning} para seleção de \textit{checkpoints}, ajuste de hiperparâmetros e ativação do critério de \textit{early stopping}. Não é utilizada na avaliação final reportada nos resultados;
	
	\item \textbf{Teste (15\%)}: constitui o \textit{gold standard} efetivo da avaliação comparativa. Todas as abordagens descritas na \autoref{sec:abordagens} --- léxicas, FinBERT-PT-BR e BERTimbau ajustado --- são avaliadas sobre exatamente este conjunto, e somente sobre ele, no cálculo das métricas reportadas no \autoref{Cap:Ava_Exp}.
\end{itemize}

A fixação da semente aleatória e a persistência dos identificadores de cada partição asseguram a reprodutibilidade integral do protocolo: qualquer execução subsequente do \textit{pipeline} sobre o mesmo corpus rotulado recupera os mesmos conjuntos, condição necessária para que comparações entre abordagens introduzidas em momentos distintos da pesquisa permaneçam metodologicamente consistentes.


\subsection{Limitações do Processo de Rotulagem}
\label{subsec:lim_rotulagem}

O processo de rotulagem foi conduzido por um único anotador — o próprio pesquisador —, o que inviabiliza o cálculo de métricas de concordância interanotadores (e.g., Kappa de Cohen). A subjetividade inerente à tarefa, em especial para publicações de polaridade ambígua ou contextualmente dependente, representa uma fonte de incerteza não quantificada nos resultados. Essa limitação
é discutida em maior profundidade na \autoref{sec:limitacoes}.
